% 2010 年浙江卷文科第 22 题：抛物线 C、准线、焦点 F，以及 △AA_1F、△BB_1F 的重心 G、H。
% 来源 PDF：普通高考/2010/2010浙江文.pdf，SHA-256 a73e22ea405697775b2d5cd82532fcffad0f506b849d9f199265b4297b8a754b。
% 原页：第 2 页（A3 三栏）第 3 栏上部，书页码 642。
% 重绘依据：照原页，抛物线开口向右、顶点在原点 O；准线为左侧竖直线，
%   A_1、B_1 在准线上且 A_1A、B_1B 水平；弦 AB 过焦点 F（即直线 l）；
%   连 A_1F、B_1F，标出重心 G、H（原页无圆，只标两点）。
%   题设 p=m^2，取 m=0.35（p=0.1225）并按 p/2 归一化作图，得 y^2=4x、准线 x=-1、
%   F(1,0)、A(1.99,2.82)、B(0.50,-1.42)、G(0.66,1.88)、H(0.17,-0.95)。
\begin{tikzpicture}[scale=1.05, semithick, every node/.style={inner sep=1.2pt, font=\small}]

  % 抛物线 y^2 = 4x，即 x = y^2/4
  \draw[domain=-1.55:3.0, samples=160, smooth, variable=\y]
    plot ({(\y*\y)/4}, \y);

  % 坐标轴
  \draw[->, >=stealth] (-1.3,0) -- (3.9,0);
  \draw[->, >=stealth] (0,-1.6) -- (0,3.05);
  \node[below right] at (3.9,0) {$x$};
  \node[left] at (0,3.05) {$y$};

  % 准线
  \draw (-1,-1.5) -- (-1,2.95);

  % 关键点
  \coordinate (A1) at (-1,2.82);
  \coordinate (B1) at (-1,-1.42);
  \coordinate (A)  at (1.99,2.82);
  \coordinate (B)  at (0.50,-1.42);
  \coordinate (F)  at (1,0);
  \coordinate (G)  at (0.663,1.88);
  \coordinate (H)  at (0.167,-0.947);

  % 水平投影线与连线
  \draw (A1) -- (A);
  \draw (B1) -- (B);
  \draw (A1) -- (F);
  \draw (B1) -- (F);
  \draw (A) -- (B);

  \fill (G) circle (1.7pt);
  \fill (H) circle (1.7pt);

  \node[left] at (A1) {$A_1$};
  \node[left] at (B1) {$B_1$};
  \node[above left] at (A) {$A$};
  \node[below right] at (B) {$B$};
  \node[below right] at (F) {$F$};
  \node[below left] at (0,0) {$O$};
  \node[left] at (G) {$G$};
  \node[right] at ($(H)+(0.06,0.09)$) {$H$};
\end{tikzpicture}
